\chapter{METODE PENELITIAN}

\section{Deskripsi Umum dan Status Protokol}

Penelitian ini dirancang sebagai evaluasi empiris terkontrol atas dua prosedur pengelolaan \textit{deployment}: prosedur imperatif terotomasi berbasis \texttt{kubectl} dan GitOps berbasis ArgoCD. Pertanyaan yang dituju ialah bagaimana kedua prosedur tersebut berbeda pada konsistensi konfigurasi, waktu deteksi, ketersediaan bukti penyimpangan, pemulihan konfigurasi, dan keterlacakan untuk skenario yang ditetapkan. Estimand dibatasi pada selisih hasil antarprosedur dalam konteks beban uji multi-service InvenioRDM dan apparatus eksperimen yang kelak diverifikasi; estimand ini bukan efek universal seluruh paradigma imperatif atau GitOps.

Naskah ini memuat rancangan protokol, bukan laporan pelaksanaan. Protokol masih memerlukan tinjauan pembimbing, dan apparatus eksperimen yang tersedia belum divalidasi untuk mengimplementasikan seluruh definisi operasional di bawah. Versi perangkat lunak, kesiapan klaster, inventaris sumber daya formal, dan data eksperimen belum ditetapkan atau diverifikasi dalam naskah ini. Karena itu, uraian rencana tidak boleh dibaca sebagai bukti bahwa eksperimen telah dilaksanakan.

\section{Rancangan Perbandingan}

Pengendalian kondisi awal, pengulangan, pencatatan penyimpangan protokol, dan pembatasan inferensi mengikuti prinsip evaluasi empiris perangkat lunak \citep{wohlin2012experimentation}. Karena pengujian memakai lingkungan yang dibangun secara sengaja pada satu klaster, rancangan tidak dianggap sebagai penugasan acak klaster dari suatu populasi \citep{cook1979quasi}.

Unit analisis adalah satu injeksi penyimpangan yang diterima pada satu lingkungan dalam satu pengulangan. Lingkungan A menggunakan otomasi imperatif tetap: pemeriksa berkala membandingkan keadaan aktual dengan inventaris dan nilai yang diinginkan, menyimpan perbedaan yang teramati, lalu menjalankan prosedur koreksi \texttt{kubectl} yang telah ditentukan. Lingkungan B menggunakan ArgoCD dengan rekonsiliasi otomatis terhadap deklarasi Git. Tidak ada operator yang menunggu, mengetik, atau memilih tindakan korektif selama pengukuran; waktu respons manusia maupun tingkat keahlian operator bukan variabel penelitian.

Kedua lingkungan memakai beban uji, nilai dasar, definisi bidang yang diperiksa, serta prosedur injeksi dan pengukuran yang setara. Injeksi dilakukan satu per satu, bukan serentak, dan urutan lingkungan serta skenario diacak dalam blok pengujian setelah kondisi dasar dipulihkan. Pengendali koreksi hanya menerima keadaan aktual dan deklarasi yang diperlukan untuk membandingkannya; identitas skenario tidak diteruskan sebagai label, termasuk pada skenario Secret. Rancangan ini membandingkan dua implementasi operasional yang telah ditentukan, bukan mengisolasi seluruh faktor yang mungkin membedakan setiap produk atau organisasi.

\begin{figure}[H]
  \centering
  \begin{tikzpicture}[
    box/.style={draw, rounded corners, align=center, text width=2.6cm, minimum height=10mm, font=\small},
    flow/.style={-{Stealth[length=2mm]}, thick},
    node distance=9mm and 5mm
  ]
    \node[box] (inject) {Injeksi\\terkontrol};
    \node[box, left=10mm of inject] (kubectl) {A: otomasi\\\texttt{kubectl}};
    \node[box, right=10mm of inject] (argocd) {B: ArgoCD dan\\rekonsiliasi Git};
    \node[box, below=of kubectl] (worka) {Beban uji\\multi-service A};
    \node[box, below=of argocd] (workb) {Beban uji\\multi-service B};
    \node[box, below=30mm of inject, text width=4.5cm] (measure) {Evaluator bersama\\R1: kesetaraan konfigurasi\\Kesehatan: laporan terpisah\\R2: waktu kejadian\\R3: bukti audit};
    \draw[flow] (inject.south) -- (worka.east);
    \draw[flow] (inject.south) -- (workb.west);
    \draw[flow] (kubectl) -- (worka);
    \draw[flow] (argocd) -- (workb);
    \draw[flow] (worka.south) |- (measure.west);
    \draw[flow] (workb.south) |- (measure.east);
  \end{tikzpicture}
  \caption{Rancangan perbandingan dua prosedur otomatis dengan evaluator dan definisi metrik yang sama.}
  \label{fig:arsitektur-eksperimen}
\end{figure}

Gambar~\ref{fig:arsitektur-eksperimen} menunjukkan bahwa injeksi, mekanisme pengelolaan, dan evaluator merupakan bagian yang dibedakan secara eksplisit. Hasil hanya akan ditafsirkan untuk konfigurasi apparatus yang benar-benar dicatat saat eksperimen.

\section{Alur Penelitian}

Tahapan yang direncanakan meliputi peninjauan protokol, validasi instrumen melalui uji pendahuluan, pengujian formal, dan analisis per skenario. Uji pendahuluan menggunakan surrogat sederhana S1 dan S2 hanya untuk memeriksa kelayakan injeksi, pencatatan waktu, pemeriksa kesetaraan, dan prosedur reset. Data S1/S2 tidak masuk ke sampel formal dan tidak digunakan untuk menyimpulkan perilaku InvenioRDM.

\begin{figure}[H]
  \centering
  \begin{tikzpicture}[
    process-box/.style={draw, rounded corners, align=center, text width=4cm, minimum height=10mm, font=\small},
    flow/.style={-{Stealth[length=2mm]}, thick},
    node distance=5mm and 10mm
  ]
    \node[process-box] (review) {Tinjauan pembimbing\\dan pembakuan protokol};
    \node[process-box, below=of review] (pilot) {Pilot S1/S2\\validasi instrumen saja};
    \node[process-box, below=of pilot] (ready) {Gerbang kesiapan\\beban uji formal};
    \node[process-box, below=of ready] (inject) {Kolektor aktif\\Satu injeksi\\Verifikasi penerimaan};
    \node[process-box, below=of inject] (measure) {Pengukuran otomatis\\hingga pulih atau batas waktu};
    \node[process-box, right=of measure] (reset) {Reset dan verifikasi\\kondisi dasar};
    \node[process-box, below=of measure] (analyze) {Analisis terpisah\\per skenario dan metrik};
    \draw[flow] (review) -- (pilot);
    \draw[flow] (pilot) -- (ready);
    \draw[flow] (ready) -- (inject);
    \draw[flow] (inject) -- (measure);
    \draw[flow] (measure) -- (reset);
    \draw[flow] (reset.north) |- node[pos=0.75, above, font=\small] {ulang} (inject.east);
    \draw[flow] (measure) -- node[right, font=\small] {selesai} (analyze);
  \end{tikzpicture}
  \caption{Bagan alir protokol, dengan pilot dipisahkan dari pengulangan formal.}
  \label{fig:bagan-alir-eksperimen}
\end{figure}

Gambar~\ref{fig:bagan-alir-eksperimen} menempatkan validasi sebelum pengumpulan data formal dan mensyaratkan pemulihan kondisi dasar di antara pengulangan. Tidak ada jeda buatan setelah injeksi untuk menunggu respons manusia.

\section{Beban Uji dan Lingkungan}

Pengujian formal direncanakan menggunakan beban uji komposit multi-service InvenioRDM dengan sasaran sedikitnya 16 komponen aplikasi dan infrastruktur pendukung sesuai ruang lingkup proposal. Sasaran tersebut bukan jumlah komponen yang telah diverifikasi. Inventaris formal mencakup seluruh sumber daya dalam manifest yang disetujui untuk kedua lingkungan; daftar komponen, batas pengelolaan, dan versinya akan dicatat dari artefak yang benar-benar digunakan. Komponen infrastruktur bersama harus dibedakan dari sumber daya yang menjadi sasaran injeksi agar isolasi kedua lingkungan dapat diperiksa. S1/S2 merupakan surrogat pilot, bukan pengganti beban formal, dan tidak boleh digabungkan ke hasil atau analisis formal.

Rancangan membatasi pengujian pada satu klaster Kubernetes dan satu implementasi GitOps, ArgoCD. Spesifikasi serta versi aktual klaster, node, komponen beban uji, dan ArgoCD harus diverifikasi dan dicatat sebelum pengujian; nilai yang pernah disebut sebagai sasaran pada draf sebelumnya belum menjadi fakta eksperimen. Dua lingkungan perlu memiliki inventaris dan nilai dasar yang sebanding serta terisolasi agar koreksi satu lingkungan tidak mengubah lingkungan lain.

\subsection{Gerbang Kesiapan}

Pengulangan formal hanya dimulai setelah seluruh prasyarat berikut diperiksa dan dicatat: (1) manifest dan inventaris sumber daya formal telah dibekukan; (2) kedua lingkungan berada pada kesetaraan nilai yang ditetapkan dan kesehatan layanan telah memenuhi kriteria kesiapan; (3) auto-sync, self-heal, cakupan aplikasi, dan kebijakan prune ArgoCD telah diverifikasi sesuai rancangan; (4) pemeriksa imperatif, pengumpul bukti, serta pembanding bidang telah lulus uji pilot; (5) jam pencatat waktu tersinkronisasi; dan (6) tidak ada pengujian atau perubahan lain yang berlangsung pada lingkup sumber daya yang diamati. Ambang kesehatan, batas beban, interval ArgoCD, serta batas waktu pengamatan akan dibakukan setelah pilot dan tinjauan pembimbing; naskah ini tidak menyatakan nilai aktualnya.

Jika prasyarat gagal, injeksi tidak dimulai. Gangguan setelah injeksi dicatat sebagai kejadian pengujian, bukan dihapus tanpa jejak. Kriteria penghentian, nilai batas waktu, dan aturan pengulangan harus ditetapkan sebelum data formal dikumpulkan.

\section{Intervensi dan Skenario}

Intervensi adalah satu perubahan terkontrol terhadap keadaan sumber daya dalam satu pengulangan. Skenario A--E dan G menyasar sumber daya yang termasuk inventaris aplikasi; skenario F menguji batas visibilitas untuk sumber daya di luar cakupan aplikasi dan dianalisis terpisah.

\begin{table}[H]
  \centering
  \small
  \caption{Skenario intervensi yang direncanakan}
  \label{tab:skenario-intervensi}
  \begin{tabularx}{\textwidth}{|c|>{\raggedright\arraybackslash}p{3.3cm}|X|}
    \hline
    Skenario & Intervensi & Batas interpretasi \\
    \hline
    A & Mengubah jumlah replika Deployment terkelola & Perubahan terhadap nilai dasar pada inventaris aplikasi \\
    \hline
    B & Mengubah referensi image Deployment terkelola & Perubahan bidang yang telah ditetapkan dalam skema pengukuran \\
    \hline
    C & Mengubah bidang ConfigMap terkelola & Tidak mengubah deklarasi sumber kebenaran \\
    \hline
    D & Menghapus sumber daya terkelola & Pemulihan dinilai terhadap inventaris dan nilai dasar \\
    \hline
    E & Mengubah bidang sumber daya terkelola melalui patch & Hanya bidang yang telah ditetapkan yang dinilai \\
    \hline
    F & Membuat atau menerapkan sumber daya pada namespace di luar cakupan aplikasi & Kasus batas visibilitas sumber daya tak terkelola; bukan keberhasilan atau kegagalan self-heal yang diasumsikan \\
    \hline
    G & Mengubah data Secret terkelola di luar deklarasi & Perbandingan dilakukan tanpa menyimpan atau menampilkan nilai Secret \\
    \hline
  \end{tabularx}
\end{table}

Pada skenario F, sumber daya di luar inventaris aplikasi tidak dianggap sebagai drift terkelola yang wajib dipulihkan ArgoCD. Pengamatan dibatasi pada apakah batas cakupan dan pencatatan inventaris membuat keberadaan sumber daya tersebut terlihat. Hasil F dilaporkan tersendiri dan tidak digabungkan ke estimasi utama A--E dan G.

\section{Definisi Variabel dan Metrik}

\subsection{R1: Kesetaraan Konfigurasi dan Kesehatan}

Kesetaraan konfigurasi dinilai menggunakan aturan yang sama pada kedua lingkungan. Sebelum eksperimen, skema bidang yang diperiksa dan inventaris sumber daya yang diharapkan dibekukan. Bidang dinamis yang tidak relevan bagi deklarasi dikecualikan secara eksplisit. Untuk setiap sumber daya yang diharapkan, sumber daya dinyatakan cocok jika keberadaannya sesuai inventaris dan semua bidang terpilih sama dengan nilai yang diinginkan.

\[
  R1_{eq} = \frac{\text{jumlah sumber daya terkelola yang cocok}}{\text{jumlah sumber daya dalam inventaris yang diharapkan}} \times 100\%.
\]

Evaluator menghitung $R1_{eq}$ dari pembacaan keadaan aktual dan inventaris yang sama bentuknya untuk kedua lengan; status \textit{Synced} ArgoCD bukan pengganti aturan kesetaraan tersebut. Nilai dicatat pada kondisi dasar, pengamatan pertama setelah injeksi, dan setelah prosedur koreksi. Jika koreksi mendahului pengamatan pascainjeksi, kondisi sesaat tersebut ditandai tidak teramati; kondisi drift tidak direkonstruksi seolah-olah telah diukur. Nilai Secret dibandingkan secara aman di memori; nilai mentah maupun hash yang dapat digunakan untuk menebak isinya tidak ditulis ke log, gambar, atau laporan.

Kesehatan layanan dilaporkan sebagai hasil terpisah berdasarkan kriteria kesiapan yang ditetapkan sebelum eksperimen, misalnya kondisi kesiapan Pod dan status layanan yang memang tersedia pada workload. Kesehatan tidak dimasukkan ke pembilang kesetaraan konfigurasi. Sumber daya ekstra di luar inventaris, termasuk F, dilaporkan terpisah dan tidak mengubah penyebut R1.

\subsection{R2: Deteksi, Ketersediaan Bukti, dan Pemulihan}

Semua waktu berasal dari satu kolektor pada host pencatat. Setiap kejadian diberi waktu jam monotonik dalam detik dengan pecahan desimal serta cap waktu UTC untuk audit; durasi dihitung dari jam monotonik, bukan dari selisih jam dinding. Pengamatan diaktifkan sebelum injeksi dan berlangsung pada interval nominal satu detik. Cap waktu adalah waktu pengamatan pertama, bukan waktu kejadian yang diinterpolasi; interval aktual dan jeda polling dicatat. Interval ini merupakan resolusi pengamatan yang direncanakan dan harus divalidasi pada pilot. Interval pemeriksa imperatif dan jadwal rekonsiliasi ArgoCD dicatat tersendiri dari interval evaluator, dibekukan sebelum pengujian formal, dan tidak disesuaikan setelah hasil terlihat.

\begin{itemize}
  \item $t_0$: waktu kolektor mulai mengirim permintaan injeksi ke Kubernetes API. Waktu penerimaan respons API juga dicatat sebagai $t_a$, sehingga $t_a-t_0$ mendokumentasikan latensi permintaan. Injeksi diterima hanya jika operasi API berhasil dan objek respons atau pembacaan balik memverifikasi perubahan bidang sasaran. Durasi sejak $t_0$ mencakup latensi permintaan, bukan semata-mata waktu respons setelah perubahan tersimpan. Injeksi gagal atau tidak terverifikasi dicatat sebagai kegagalan injeksi, bukan sebagai waktu deteksi nol.
  \item $t_d$: waktu pertama mekanisme pengelolaan memberi sinyal deteksi berdasarkan perbedaan dari nilai yang diinginkan. Pada A, sinyal berasal dari pemeriksa imperatif; pada B, dari status perbedaan ArgoCD yang relevan. Tidak ada sinyal skenario tersembunyi yang diberikan kepada mekanisme tersebut.
  \item $t_e$: waktu pertama bukti dari mekanisme pengelolaan tersimpan dan mengidentifikasi sumber daya serta bidang yang berbeda terhadap acuan. Bukti harus dapat diperiksa ulang tanpa memakai label injeksi yang tersembunyi; untuk Secret cukup dicatat identitas bidang dan hasil ketidakcocokan, bukan isinya. Kriteria ini hanya mengukur bukti deviasi, bukan kelengkapan seluruh unsur R3 atau identifikasi penyebab. $t_e$ dapat sama dengan $t_d$, tetapi keduanya tetap dicatat terpisah. Pengumpulan bukti tidak menunda koreksi pada salah satu lingkungan.
  \item $t_r$: waktu pengamatan pertama ketika inventaris dan seluruh bidang R1 kembali cocok dengan nilai yang diinginkan, dengan kecocokan dikonfirmasi pada pembacaan berikutnya. Waktu yang dilaporkan tetap waktu pengamatan pertama; konfirmasi hanya mencegah keadaan sesaat dinyatakan pulih. Kesehatan layanan dilaporkan terpisah.
\end{itemize}

Durasi utama adalah waktu deteksi $t_d-t_0$, waktu ketersediaan bukti $t_e-t_0$, dan waktu pemulihan konfigurasi penuh $t_r-t_0$. Selang dari deteksi ke bukti ($t_e-t_d$) dapat ditampilkan sebagai rincian tambahan. $t_r$ berarti pemulihan kesetaraan inventaris dan bidang yang diukur, bukan klaim bahwa seluruh layanan sehat atau tidak mengalami dampak lain.

Sinyal deteksi A berasal dari pemeriksaan otomatis terhadap keadaan aktual; setelah sinyal muncul, prosedur koreksi tetap menerapkan deklarasi yang diinginkan menggunakan tindakan \texttt{kubectl} yang telah ditetapkan. B menggunakan rekonsiliasi ArgoCD. Kedua prosedur berjalan otomatis tanpa waktu reaksi manusia dan tanpa jeda buatan. Keputusan koreksi menggunakan hasil perbandingan umum, bukan label A--G yang disembunyikan dari responden.

\subsection{R3: Rubrik Bukti Keterlacakan}

R3 dinilai deskriptif untuk setiap pengulangan dengan rubrik yang sama, bukan diasumsikan lengkap karena salah satu alat digunakan. Penilai mencatat ada, sebagian, atau tidak ada bukti untuk unsur berikut.

\begin{table}[H]
  \centering
  \small
  \caption{Rubrik penilaian bukti R3}
  \label{tab:rubrik-r3}
  \begin{tabularx}{\textwidth}{|>{\raggedright\arraybackslash}p{2.5cm}|X|>{\raggedright\arraybackslash}p{4cm}|}
    \hline
    Unsur & Kriteria bukti & Contoh sumber yang diperiksa \\
    \hline
    Pelaku & Identitas aktor atau akun layanan yang melakukan perubahan dapat ditelusuri & Log audit API Kubernetes; log eksekusi otomatis \\
    \hline
    Objek dan perubahan & Objek serta bidang yang berubah dapat ditentukan tanpa menampilkan data sensitif & Diff ArgoCD yang telah disanitasi; catatan perbandingan imperatif \\
    \hline
    Waktu dan urutan & Waktu serta urutan injeksi, deteksi, dan pemulihan dapat direkonstruksi & Cap waktu kolektor; log audit; metadata Git \\
    \hline
    Tindakan dan hasil & Tindakan koreksi dan hasil verifikasinya dapat dikaitkan dengan pengulangan & Log koreksi; riwayat sinkronisasi; pembacaan verifikasi \\
    \hline
  \end{tabularx}
\end{table}

Rubrik ini mengukur ketersediaan bukti perubahan dan tindakan. Rubrik tidak meminta penilai menyimpulkan penyebab kausal atau akar masalah dari sebuah diff. Data audit dapat memuat identitas sensitif; pelaporan hanya menggunakan informasi minimum yang telah disanitasi.

\section{Protokol Pengulangan dan Penanganan Kegagalan}

Satu pengulangan mengikuti urutan berikut:
\begin{enumerate}
  \item Periksa gerbang kesiapan dan simpan bukti kondisi dasar, inventaris, kesehatan, serta status instrumentasi.
  \item Aktifkan koleksi sinyal deteksi, bukti, kesetaraan, dan kesehatan sebelum mengirim injeksi; pastikan mekanisme pengelolaan juga telah aktif. Jalankan hanya satu skenario pada satu lingkungan. Catat $t_0$, respons API pada $t_a$, dan verifikasi penerimaan injeksi tanpa menambahkan label atau identitas skenario ke sumber daya yang diamati.
  \item Pertahankan koleksi berkala satu detik tanpa menyisipkan waktu tunggu agar manusia bereaksi atau agar drift bertahan selama durasi tertentu. Jika perubahan diterima tetapi sinyal singkat tidak teramati, laporkan keterbatasan observasi tanpa menganggap deteksi terjadi pada waktu nol.
  \item Lanjutkan sampai pemulihan terkonfirmasi atau batas waktu pengamatan yang telah dibakukan tercapai. Simpan semua kejadian, jeda polling, kesalahan API, serta status penghentian.
  \item Pulihkan kondisi dasar, verifikasi kesetaraan inventaris dan kesehatan, lalu baru mulai pengulangan berikutnya.
\end{enumerate}

Batas waktu maksimum pengamatan harus ditentukan dan dibekukan setelah pilot serta tinjauan pembimbing, sebelum pengumpulan formal. Tidak ada nilai batas waktu yang diklaim telah diuji dalam draf ini. Jika sinyal deteksi, bukti, atau pemulihan belum terjadi saat batas itu tercapai, hasil masing-masing ditandai tersensor kanan pada batas tersebut; waktu hilang tidak diganti nol dan pengulangan tidak dikeluarkan hanya karena pemulihannya gagal. Kegagalan injeksi, kegagalan instrumentasi, gangguan infrastruktur, dan tidak terjadinya deteksi merupakan kategori berbeda. Setiap kategori, jumlah pengulangan terdampak, alasan, serta aturan pengulangan yang telah disetujui dilaporkan.

Run dengan respons injeksi sukses tetapi tanpa bukti bahwa perubahan bidang sasaran pernah diterima tidak menjadi hasil nol; run ditandai sebagai injeksi tidak valid dan dicatat. Jika perubahan diterima tetapi mekanisme tidak mendeteksinya sampai batas pengamatan, itu merupakan hasil deteksi tersensor, bukan kegagalan instrumen. Jika kolektor atau sumber data gagal sehingga status tidak dapat ditentukan, run ditandai kegagalan instrumentasi dan tidak ditafsirkan sebagai deteksi terlewat. Skenario F tetap diklasifikasikan sebagai pengamatan batas visibilitas dan tidak dinilai dengan asumsi pemulihan otomatis.

\section{Pengulangan dan Analisis Data}

Target rancangan adalah $n=5$ pengulangan per skenario dan per lingkungan untuk R1 serta $n=10$ untuk R2. Angka tersebut merupakan target deskriptif yang direncanakan, bukan hasil analisis daya dan bukan jaminan kecukupan daya untuk setiap skenario. Pengulangan R3 mengikuti pengumpulan bukti pada run yang sama dan dianalisis secara deskriptif. S1/S2 tidak dihitung dalam target ini. Skenario F dilaporkan tersendiri dan tidak dimasukkan ke ringkasan utama skenario drift terkelola A--E dan G.

Hasil disajikan per skenario dan per lingkungan: jumlah run direncanakan, jumlah injeksi valid, jumlah kejadian teramati dan tersensor, ringkasan distribusi waktu, proporsi pemulihan pada batas pengamatan, R1, kesehatan terpisah, serta rubrik R3. Nilai mentah yang hilang tidak diimputasi sebagai nol. Observasi yang tersensor tidak diam-diam dibuang agar hanya run yang berhasil pulih yang dibandingkan. Jika penyensoran terjadi, ringkasan dan uji yang tidak memperhitungkan penyensoran tidak ditafsirkan sebagai perbandingan lengkap; analisis komparatif untuk hasil tersebut harus ditetapkan sebelum melihat hasil, atau dilaporkan secara deskriptif tanpa klaim inferensial.

Untuk hasil waktu yang teramati lengkap, rencana analisis eksploratif menggunakan uji Mann--Whitney U dua sisi apabila asumsi independensi pengulangan terpenuhi. Cliff's $\delta$ ditetapkan sebagai ukuran efek utama agar pilihan ukuran efek tidak bergantung pada arah hasil yang terlihat. Distribusi, median, rentang, dan ukuran sampel yang benar-benar tersedia tetap dilaporkan. Jika nilai $p$ digunakan, ambang nominalnya adalah $\alpha=0{,}05$ dengan penyesuaian Holm untuk keluarga perbandingan waktu yang dibakukan sebelum pengumpulan formal. Definisi keluarga uji dan pemenuhan independensi harus ditinjau bersama pembimbing sebelum protokol dibekukan; tidak ada klaim konfirmatoris dari draf ini. Penggabungan lintas skenario tidak menjadi analisis utama karena bentuk dan dampak intervensinya berbeda. R1, kesehatan, R3, dan kasus F dilaporkan deskriptif; tidak ada hasil signifikansi atau besar efek yang dinyatakan sebelum data dikumpulkan.

\section{Ancaman terhadap Validitas dan Keterbatasan Rancangan}

Validitas internal dapat terancam oleh perbedaan beban klaster dari waktu ke waktu, perubahan konfigurasi selain intervensi, ketidaksetaraan inventaris, serta perbedaan jadwal rekonsiliasi. Isolasi namespace, nilai dasar yang dibekukan, urutan pengujian yang diacak, pencatatan keadaan, dan gerbang kesiapan dirancang untuk mengurangi ancaman tersebut, tetapi tidak menjamin bahwa satu-satunya perbedaan adalah paradigma abstrak. Perbandingan tetap mengikutsertakan perbedaan implementasi yang melekat pada pemeriksa/korektor \texttt{kubectl} dan ArgoCD.

Validitas konstruk dibatasi oleh tujuh skenario yang dipilih, definisi bidang yang dibandingkan, interval polling, batas waktu, kriteria kesehatan, dan rubrik bukti. Pengamatan polling memberikan waktu deteksi pada resolusi instrumen, bukan waktu kejadian kontinu yang persis. Skenario F menguji batas inventaris dan tidak dapat digeneralisasi sebagai pengukuran self-healing untuk sumber daya terkelola.

Validitas eksternal dibatasi pada satu klaster, satu beban uji formal multi-service InvenioRDM, satu alat GitOps, serta konfigurasi dan versi yang kelak benar-benar digunakan. Pilot S1/S2 tidak memperluas cakupan inferensi. Tidak ada kesimpulan yang dapat digeneralisasi ke semua aplikasi, klaster, atau alat GitOps dari rancangan ini saja.

Reliabilitas bergantung pada konsistensi kolektor, sinkronisasi jam, cakupan log, dan kemampuan mengulang reset. Sebelum data formal dikumpulkan, keberhasilan pencatatan, pembandingan Secret tanpa kebocoran, deteksi kegagalan, serta aturan pemulihan dan penyensoran harus diuji. Sampai gerbang tersebut dipenuhi, hasil eksperimen belum tersedia dan protokol tetap berstatus draf.
